\section{Cartesian motion of the complete real-order zero arc}
\label{rwcart:sec:motion}
The normalized radial motion can have different signs above the critical
line, but its Cartesian abscissa has one direction throughout the entire
nonnegative-outer-order domain.

\begin{theorem}[Strict outward abscissa motion]
\label{rwcart:thm:motion}
For every $a\ge0,b>0$, put $s=\rho^2$ and
$X_{a,b}(s)=\rho\cos\theta_{a,b}(\rho)$. The unique upper-arc zero
has a real-analytic abscissa on $(0,1]$, extending real analytically
through zero, with
\[
 0<X_{a,b}(s)<s,\qquad X_{a,b}'(s)>0\quad(0<s\le1).
\]
This is an abscissa theorem; it does not assert one direction for
$\cos\theta/\rho$ in every supercritical parameter pair.
\end{theorem}
\begin{proof}
For $a>0,a+b>1$, let $\nu$ be the finite signed measure of
Theorem~\ref{subunit:thm:domain}. Set
\[
 G(X,s)=\int_0^1\frac{d\nu(u)}{D(u)},\qquad
 D(u)=1-2Xu+su^2.
\]
The angular root satisfies $G=0$ and $0<X<s\le1$.
At that root the measure $d\tau=d\nu/D$ has zero total mass and
the same ordered signs. The sign test gives
\begin{align*}
 G_X&=2\int\frac uD\,d\tau>0,
 &\left(\frac uD\right)'&=\frac{1-su^2}{D^2}>0,\\
 G_s&=-\int\frac{u^2}D\,d\tau<0,
 &\left(\frac{u^2}D\right)'&=\frac{2u(1-Xu)}{D^2}>0.
\end{align*}
Thus $X'=-G_s/G_X>0$. At $(X,s)=(0,0)$, $G=0$ and
$G_X=2c_2>0$, proving analytic extension by the implicit-function theorem.

For $a+b\le1$ and on the full axis, the right-endpoint factorization
gives $X=s\eta(\sqrt s)$, with $\eta>0$ and
$\partial_s\eta>0$ by Theorem~\ref{subpick:thm:radial}.
Its positive-measure equation has derivative $G_m=\omega([0,1])>0$
at $s=0$, so $m$ and $\eta=(1+m)/2$ are analytic there.
This proves the same derivative conclusion. In every case the root
is away from the cut at $s=1$, and all denominators are locally
uniformly nonzero; the analytic implicit argument also applies there.
\end{proof}

\subsection{The opposite endpoint mechanism in the large-inner-order limit}
For fixed $a>0$, the disk coefficients tend as $b\to\infty$ to
$c_1=0$, $c_n=n^{-a}$ for $n\ge2$. The limit function is
$F_{a,\infty}(z)=\Li_a(z)-z$. Its signed measure is the positive
Gamma moment law minus a unit atom at zero. Thus its normalized-radius
motion has the opposite direction from the right-endpoint case.

\begin{theorem}[Strict decrease for the left-endpoint limit]
\label{rwcart:thm:left-atom}
For every $a>0$, the unique upper-arc zero of
$\Imag F_{a,\infty}$ satisfies
\[
 \arccos(\rho/2)<\theta_{a,\infty}(\rho)<\pi/2,
 \qquad 0<\eta_{a,\infty}(\rho)<\tfrac12,
 \qquad \eta_{a,\infty}'(\rho)<0
\]
for $0<\rho\le1$.
\end{theorem}
\begin{proof}
Let $U=e^{-Y}$, $Y\sim\mathrm{Gamma}(a,1)$, and let
$d\omega(u)=u\,d\Pr(U\le u)$. The measure has strictly positive
density on $(0,1)$ and mass $2^{-a}$. Moment expansion gives
\[
 F_{a,\infty}(z)=z^2\int_0^1\frac{d\omega(u)}{1-zu}.
\]
Set $x=\rho^2$, $m=2\cos\theta/\rho$, and
\[
 D(u)=1+x(u^2-mu),\qquad
 G(m,x)=\int_0^1\frac{m-u}{D(u)}\,d\omega(u).
\]
Direct rationalization shows that the imaginary part has the sign of
$G$. We have $G(0,x)<0<G(1,x)$, and
$G_m=\int D^{-2}\,d\omega>0$. Outside $0<m<1$ the numerator has
a fixed sign. Thus the zero is unique and simple on the full upper arc.
At the root, differentiation gives
\[
 G_x=\int\frac{u(u-m)^2}{D(u)^2}\,d\omega(u)>0,
 \qquad
 \frac{dm}{dx}=-\frac{\int u(u-m)^2D(u)^{-2}\,d\omega(u)}
                       {\int D(u)^{-2}\,d\omega(u)}<0.
\]
All denominators are uniformly positive at $0<m<1$, $0\le x\le1$.
Hence $\eta'(\rho)=\rho m'(\rho^2)<0$, proving the claims.
The theorem concerns the exact limit family, not an unproved transfer
of its radial sign to every finite inner order.
\end{proof}
